\documentclass{article} % For LaTeX2e
\usepackage{iclr2027_conference,times}

% Optional math commands from https://github.com/goodfeli/dlbook_notation.
\input{math_commands.tex}

\usepackage{graphicx}
\usepackage{booktabs}
\usepackage{multirow}
\usepackage{array}
\usepackage{amssymb}
\usepackage{wrapfig}
\usepackage{placeins}
\usepackage{float}
\usepackage{flafter}
\usepackage{longtable}
\usepackage{hyperref}
\hypersetup{hidelinks}
\makeatletter
\setlength{\@fptop}{0pt}
\setlength{\@fpsep}{10pt}
\setlength{\@fpbot}{0pt plus 1fil}
\makeatother
\usepackage{xurl}
\providecommand{\doi}[1]{\href{https://doi.org/#1}{doi:#1}}
\usepackage[bottom]{footmisc}
\setcounter{topnumber}{3}
\setcounter{bottomnumber}{2}
\setcounter{totalnumber}{4}
\renewcommand{\topfraction}{0.95}
\renewcommand{\bottomfraction}{0.4}
\renewcommand{\textfraction}{0.05}
\renewcommand{\floatpagefraction}{0.8}
\setlength{\textfloatsep}{8pt plus 2pt minus 4pt}
\setlength{\intextsep}{8pt plus 2pt minus 3pt}

\title{Global Grammar Meets Local Constraint: A Training-Free Readout for Protein Mutation Prediction}

% Authors must not appear in the submitted version. They should be hidden
% as long as the \iclrfinalcopy macro remains commented out below.
% Non-anonymous submissions will be rejected without review.

\author{Anonymous authors\\
Paper under double-blind review
}

\newcommand{\fix}{\marginpar{FIX}}
\newcommand{\new}{\marginpar{NEW}}

%\iclrfinalcopy % Uncomment for camera-ready version, but NOT for submission.
\begin{document}

\maketitle

\begin{abstract}
Pretrained protein models have advanced mutation-effect prediction, yet their frozen scores primarily express a corpus-level amino-acid grammar rather than constraints specific to the queried protein family. Across $37$ checkpoints, we find that protein-wide native preferences track corpus abundance with Pearson correlation 0.99, whereas homolog-derived preferences vary across sites and become most informative when they redirect the frozen ranking. This separation is strongest where overlap drops to only 19.7% at a similarity threshold of 0.3, making family evidence a complementary inference signal. We introduce \textsc{REM2}, an architecture- and task-agnostic, training-free inference-time readout that adaptively integrates homolog evidence with frozen model predictions through entropy weighting, coherence-guided calibration, and structure-aware uncertainty correction. A closed-form decomposition isolates the contribution of each operation. Across masked language models, autoregressive models, structure-aware encoders, and inverse-folding models, we evaluate $59$ configurations on $1{,}211$ assays spanning high-throughput DMS from ProteinGym, high-precision low-throughput measurements from VenusMutHub, and VenusViroHub as an independent viral test. \textsc{REM2} improves every configuration across all three benchmarks and exhibits broad gains under five complementary metrics. Separation, domain, and functional analyses connect global--local disagreement to biologically organized changes in mutation ranking. Built with \textsc{REM2}, \textsc{VenusREM2} reaches an official ProteinGym Average Spearman of $0.556$, exceeds the highest public score by $0.038$, and ranks first across all five functional categories.
\end{abstract}

\section{Introduction}

\begin{wrapfigure}{r}{0.49\textwidth}
\vspace{-1.0em}
\centering
\includegraphics[width=\linewidth]{figures/fig1_bias_teaser.png}
\vspace{-0.8em}
\caption{Consensus amino acid background across 37 protein model checkpoints compared with UniProtKB/Swiss-Prot composition (Pearson $r=0.99$). \textsc{ESM3} is the main departure, consistent with its multimodal pretraining.}
\label{fig:teaser}
\vspace{-1.0em}
\end{wrapfigure}

\begin{figure}[!t]
\centering
\includegraphics[width=\textwidth]{figures/fig2_motivation.pdf}
\vspace{-0.5em}
\caption{Global composition agreement coexists with local ranking divergence. \textbf{(a)} Family preferences vary across sites in a representative alignment. \textbf{(b)} Dataset-colored site profiles overlap but rarely align with both corpus abundance and the protein-level native background: $19.7\%$ for ProteinGym, $24.8\%$ for VenusMutHub, and $38.1\%$ for VenusViroHub at $r\geq0.3$. A fixed sample of $7{,}500$ sites per benchmark is shown; rates use all complete profiles.}
\label{fig:motivation}
\vspace{-0.5em}
\end{figure}

Single amino acid substitutions can reshape protein stability, catalysis, molecular recognition, and cellular fitness, making their prediction central to variant interpretation and protein engineering \citep{zhou2024mlife,notin2024proteingym}. Deep mutational scanning provides direct measurements but covers only a small fraction of natural proteins and possible variants. Computational models offer a scalable route by ranking mutations from sequence or structure without assay labels \citep{meier2021esm1v,lin2023esm2,frazer2021eve}.

Two forms of evolutionary information support this task. Protein language models trained on large sequence corpora learn regularities shared across families, while alignment models extract constraints from homologs of a particular protein \citep{rives2021esm1b,hopf2017evmutation}. Global representations provide transferable priors; homologs record the substitutions tolerated at each site. \textsc{Tranception} and \textsc{VenusREM} show that family retrieval can strengthen pretrained models at inference time \citep{notin2022tranception,tan2025venusrem}. Yet these combinations are designed around a particular backbone or mixing rule. A general method must determine what each signal contributes and when their information differs.

We expose this distinction in frozen mutation scores. Across $37$ checkpoints spanning masked, autoregressive, structure-aware, and inverse-folding models, protein-wide native profiles converge to a common amino acid ordering and track UniProtKB/Swiss-Prot abundance (Pearson $r=0.99$; Figure~\ref{fig:teaser}). This shared direction reflects a global grammar inherited from pretraining corpora. MSA profiles instead reorganize amino acid preferences at individual sites according to family conservation. Retrieval gains rise where the native and family fields diverge, showing that family history is most informative when it redirects rather than repeats the global preference.

Motivated by this observation, we present \textsc{REM2}, an architecture-agnostic readout that requires no training. \textsc{REM2} maps any backbone to an $L\times20$ amino acid score field and combines it with a family field on the same coordinates. Entropy weighting controls family evidence, a coherence gate calibrates the shared native background, and structural exposure and confidence adjust uncertain site scores. The result is a lightweight closed-form operator that transfers across sequence, sequence-structure, and inverse-folding models without retraining.

We support this formulation through mathematical decomposition, broad experiments, and biological analysis. The decomposition isolates family retrieval, native background calibration, and structural shrinkage. Across four model families, $59$ configurations, $1{,}211$ assays, and five metrics, \textsc{REM2} improves every configuration on ProteinGym, VenusMutHub, and VenusViroHub; all $295$ model--metric comparisons improve on both ProteinGym and VenusViroHub, and $293$ of $295$ improve on VenusMutHub. Analyses of separation regimes, protein domains, and functional categories trace the gains from site-level disagreement to protein phenotype. Using \textsc{REM2}, we build \mbox{\textsc{VenusREM2}}, a leaderboard model that reaches $0.556$ Average Spearman under the official ProteinGym protocol and ranks first across all five functional categories.

\begin{figure}[!t]
\begin{center}
\includegraphics[width=0.82\textwidth]{figures/fig3_framework.png}
\end{center}
\vspace{-0.6em}
\caption{\textsc{REM2} readout: set the family-field contribution relative to native field $R$, residualize the coherent native background, and shrink mutation contrasts by exposure and confidence (Equations~\ref{eq:fusion}--\ref{eq:final_field}).}
\label{fig:framework}
\vspace{-0.5em}
\end{figure}

\section{Empirical Motivation: Global and Family Signals Operate at Different Scales}
\label{sec:two_scale_hypothesis}

\paragraph{Global agreement and local rotation occur at different scales.}
Across $37$ masked, autoregressive, structure-aware, and inverse-folding checkpoints, protein-level native profiles have median pairwise correlation $0.995$ ($0.993$ across model families), and their consensus tracks log UniProtKB/Swiss-Prot abundance at Pearson $r=0.99$ (Figure~\ref{fig:teaser}). Homologs retain similar aggregate composition but can reorder amino-acid alternatives at individual sites. Across ProteinGym, VenusMutHub, and VenusViroHub, family profiles correlate with corpus abundance at medians of $0.206$, $0.249$, and $0.273$, and with the protein-level native profile at medians of $0.205$, $0.229$, and $0.332$, respectively. Only $19.7\%$, $24.8\%$, and $38.1\%$ align with both at $r\geq0.3$ (Figure~\ref{fig:motivation}).

\paragraph{The scale separation motivates a conditional readout.}
Agreement can repeat the native ordering, whereas local rotation can contribute a complementary ranking direction. We therefore hypothesize that family evidence is most useful when it redirects rather than reproduces the native field. \textsc{REM2} operationalizes this distinction by retaining native and family fields separately, adapting their evidence share, calibrating only a coherent native background, and applying structural reliability after fusion. Section~\ref{sec:results_scales} tests the separation--gain prediction under protein-level controls.

\section{\textsc{REM2}: A Shared Inference-Time Readout}
\label{sec:method}

\textsc{REM2} treats mutation prediction as calibration of two fixed evidence fields: a native field from a frozen backbone and a family field from retrieved homologs. The backbone parameters and conditioning context never change. The method instead solves three inference-time problems in closed form: how much family evidence to admit, when a protein-wide native direction should be removed, and how structural reliability should attenuate the resulting contrasts (Figure~\ref{fig:framework}).

\begingroup
\setlength{\abovedisplayskip}{2pt plus 1pt minus 1pt}
\setlength{\belowdisplayskip}{2pt plus 1pt minus 1pt}
\setlength{\abovedisplayshortskip}{0pt plus 1pt}
\setlength{\belowdisplayshortskip}{1pt plus 1pt minus 1pt}
\setlength{\jot}{1pt}

\subsection{A common score-field interface}
\label{sec:interface}

\paragraph{Common mutation coordinates isolate the readout.}
Let $\mathcal A$ be the 20 amino acids and $x$ a length-$L$ wild-type sequence. A frozen backbone exposes $R\in\mathbb{R}^{L\times20}$; with $q_{\theta,i}(a;C_i)$ its distribution at site $i$, masked, wild-type, and teacher-forced interfaces become
\begin{equation}
\begin{aligned}
R^{\mathrm{mask}}_{i,a}&=\log q_{\theta,i}(a;x_{\setminus i}), &
R^{\mathrm{wt}}_{i,a}&=\log q_{\theta,i}(a;x), &
R^{\mathrm{tf}}_{i,a}&=\log q_{\theta,i}(a;x_{<i}).
\end{aligned}
\label{eq:readouts}
\end{equation}
Required structure enters $C_i$; both $C_i$ and frozen $\theta$ are fixed within every Raw--\textsc{REM2} pair.

\begin{samepage}
Let $\mathcal I$ contain the $n$ query-mapped sites and $n_{i,v}$ be the retained-row count of tokenizer symbol $v$. With equal row weights and no pseudocounts, the aligned frequency and log-softmax family field are
\begin{equation}
f_{i,a}=\frac{n_{i,a}}{\sum_{v\in\mathcal V}n_{i,v}},
\qquad
E_{i,a}=f_{i,a}-\operatorname{logsumexp}_{a'\in\mathcal A}f_{i,a'},
\qquad i\in\mathcal I.
\label{eq:family_field}
\end{equation}
$\mathcal V$ includes padding, while $E$ retains amino-acid columns. The log-softmax is a design choice rather than a probabilistic model: it keeps the family field bounded and monotone in the column frequencies, is approximately linear in $f_{i,a}$ for the rare amino acids that carry most of the mutation signal, and places family contrasts on a scale commensurate with the log-probability native field before the scale weighting of Equation~\ref{eq:amplitude_share}. We set $E_i=R_i$ outside $\mathcal I$, making retrieval the identity on unmapped sites.
\end{samepage}
For either field $X$, the shared mutation coordinate is
\begin{equation}
\Delta^X_{i,a}=X_{i,a}-X_{i,x_i}.
\label{eq:contrasts}
\end{equation}
This contrast removes site-wise offsets before any evidence is combined.

\subsection{Adaptive retrieval as scale allocation}
\label{sec:fusion}

\paragraph{A target evidence share determines a unique mixing coefficient.}
On mapped sites, normalize the native field as
\begin{equation}
p_i(a)=\frac{\exp R_{i,a}}{\sum_{a'\in\mathcal A}\exp R_{i,a'}},
\qquad i\in\mathcal I.
\label{eq:cue_marginal}
\end{equation}
Write $[z]_0^1=\min\{1,\max\{0,z\}\}$ and $[z]_+=\max(z,0)$. Native uncertainty is summarized locally by $H_i$ and at protein level by $\bar H$:
\begin{equation}
H_i=-\sum_{a\in\mathcal A}p_i(a)\log p_i(a),
\qquad
\bar H=\left[\frac{1}{n\log 20}\sum_{i\in\mathcal I}H_i\right]_0^1.
\label{eq:normalized_entropy}
\end{equation}
Two diagnostics then measure weak native support and native--family rotation:
\begin{equation}
\begin{aligned}
\pi&=\frac{1}{n}\sum_{i\in\mathcal I}\log p_i(x_i), &
\rho_\pi&=[-\pi/\log 20]_0^1,\\
r_{R,E}&=\frac{1}{n}\sum_{i\in\mathcal I}\operatorname{corr}_{a\in\mathcal A}(R_i,E_i), &
\rho_{\mathrm{rot}}&=\frac{1-r_{R,E}}{2}.
\end{aligned}
\label{eq:field_alignment}
\end{equation}
Zero-variance correlations are set to zero. The desired family share is
\begin{equation}
\rho_{\mathrm{MSA}}=(1-\bar H)\rho_\pi+\bar H\rho_{\mathrm{rot}}.
\label{eq:pref_map}
\end{equation}
Thus confident native fields use wild-type support to set retrieval, whereas uncertain fields rely on whether family evidence changes the native direction.

Let $s_R$ and $s_E$ be the standard deviations of non-wild-type contrasts on $\mathcal I$. We parameterize the mix by a target evidence share rather than interpolating raw numerical scales, requiring the scale-weighted family share to equal $\rho_{\mathrm{MSA}}$:
\begin{equation}
\rho_{\mathrm{MSA}}
=\frac{\alpha s_E}{(1-\alpha)s_R+\alpha s_E}.
\label{eq:amplitude_share}
\end{equation}
Solving this constraint for the mixing coefficient gives
\begin{equation}
\alpha
=\frac{\rho_{\mathrm{MSA}}s_R}
{\rho_{\mathrm{MSA}}s_R+(1-\rho_{\mathrm{MSA}})s_E}.
\label{eq:native_alpha}
\end{equation}
The coefficient therefore follows from a prescribed evidence share rather than an unscaled interpolation. It gives the retrieved field and contrast
\begin{equation}
F_{i,a}=(1-\alpha)R_{i,a}+\alpha E_{i,a},
\qquad
\Delta^F_{i,a}=(1-\alpha)\Delta^R_{i,a}+\alpha\Delta^E_{i,a}.
\label{eq:fusion}
\end{equation}

\subsection{Coherence-gated calibration of the native channel}
\label{sec:bias}
\label{sec:calibration}

\paragraph{Only a locally coherent global direction is residualized.}
The protein-wide native direction is the log-mean-exp profile of $R$, standardized across amino acids:
\begin{equation}
b_a=\operatorname{logsumexp}_{i=1}^{L}R_{i,a}-\log L,
\qquad
\widehat b_a=\frac{b_a-\bar b}{s_b}.
\label{eq:background}
\end{equation}
Set $\widehat b=0$ when $s_b=0$. Its mean site-wise coherence and one-sided gate are
\begin{equation}
c=\frac{1}{L}\sum_{i=1}^{L}\operatorname{corr}_{a\in\mathcal A}(R_i,b),
\qquad
\kappa(R)=[c]_+.
\label{eq:gate}
\end{equation}
When their mean alignment is non-positive, $\kappa(R)=0$ and no global correction is imposed. Otherwise, calibration removes the protein-wide direction only from the native channel:
\begin{align}
S^{\mathrm{cal}}_{i,a}
&=(1-\alpha)\Delta^R_{i,a}+\alpha\Delta^E_{i,a}
 -(1-\alpha)\kappa(R)(\widehat b_a-\widehat b_{x_i})\notag\\
&=\Delta^F_{i,a}-(1-\alpha)\kappa(R)(\widehat b_a-\widehat b_{x_i}).
\label{eq:calibrated}
\end{align}
The factor $1-\alpha$ preserves the family contribution while calibrating the native background.

\input{table1_leaderboard.tex}

\subsection{Reliability-adjusted output and variant scoring}
\label{sec:shrinkage}
\label{sec:ranking}

\paragraph{Reliability completes the closed-form operator.}
Let $r_i$ be RSA, $\ell_i$ normalized pLDDT, and $q_i=1-\ell_i$, all in $[0,1]$. Relative to mapped-site means, define
\begin{equation}
\lambda_i
=\bigl(1-[r_i-\bar r]_+\bigr)
 \bigl(1-[q_i-\bar q]_+\bigr).
\label{eq:lambda}
\end{equation}
Only above-mean exposure or uncertainty attenuates a site, and neither factor can reverse a contrast. The complete readout is
\begin{equation}
\begin{aligned}
S_{i,a}
&=\lambda_i S^{\mathrm{cal}}_{i,a}\\
&=\lambda_i\!\left[\Delta^F_{i,a}
 -(1-\alpha)\kappa(R)(\widehat b_a-\widehat b_{x_i})\right].
\end{aligned}
\label{eq:final_field}
\end{equation}
A missing structural covariate contributes a factor of one (Appendix~\ref{app:structure_extension}). For a mutation set $\mathcal M$, \textsc{REM2} sums the resulting main effects:
\begin{equation}
S(\mathcal M)=\sum_{(i,a)\in\mathcal M}S_{i,a}.
\label{eq:aggregate}
\end{equation}
For \textsc{VenusREM2}, let $k\in\mathcal K=\{20,128,512,1024,2048,4096\}$ index the six \textsc{ProSST} structure vocabularies. We standardize each member's scores over the variant set $\mathcal V_A$ of assay $A$ and average them:
\begin{equation}
S_{\mathrm{VenusREM2}}(v)
=\frac{1}{6}\sum_{k\in\mathcal K}
\frac{S_k(v)-\mu_{k,A}}{\sigma_{k,A}}.
\label{eq:prosst_ensemble}
\end{equation}
Here $S_k(v)$ is Equation~\ref{eq:aggregate} evaluated with member $k$ and its entropy-adaptive $\alpha_k$; $\mu_{k,A}$ and $\sigma_{k,A}$ are the mean and standard deviation over $v\in\mathcal V_A$. Thus each vocabulary contributes equally despite score-scale differences.
Given $R$ and $E$, the protein-level operator costs $O(20L)$ and an $M$-substitution variant costs $O(M)$ indexed lookup; $\theta$ remains frozen.

\endgroup

\begin{figure}[!t]
\centering
\includegraphics[width=\textwidth]{figures/fig4_multimetric.pdf}
\caption{Paired Raw--\textsc{REM2} performance across five metrics on ProteinGym (top), VenusMutHub (middle), and VenusViroHub (bottom). Colors denote backbone families; points above the diagonal favor \textsc{REM2}.}
\label{fig:scatter}
\end{figure}

\section{Results}
\label{sec:benchmarks}



\subsection{Experimental setup}
\label{sec:setup}

\paragraph{Benchmarks.}
We evaluate three complementary benchmarks: ProteinGym ($217$ standardized substitution assays), VenusMutHub ($905$ assays across $527$ wild-type sequence clusters), and VenusViroHub ($89$ viral assays with zero overlap with ProteinGym). VenusViroHub contains $646{,}756$ substitutions from $19$ viruses across seven phenotype classes, extending evaluation to immune escape and cell entry. Appendix~\ref{app:virogym} details its scope and curation.

\paragraph{Baseline models.}
We form a shared panel of $59$ paired Raw--\textsc{REM2} configurations from $37$ frozen checkpoints spanning masked and autoregressive language models, structure-aware encoders, and inverse-folding models; every configuration is evaluated unchanged on all three benchmarks. Each pair fixes the checkpoint, structural inputs, and native scoring interface, isolating the inference-time readout. \textsc{VenusREM2} applies \textsc{REM2} independently to six \textsc{ProSST} structure vocabularies, then $z$-scores and averages their per-assay variant-score vectors. Appendix~\ref{app:baselines} lists the full model--readout taxonomy and implementations.

\paragraph{Metrics.}
Spearman correlation is the primary ranking endpoint; NDCG, AUC, MCC, and top-variant recall measure complementary ranking, discrimination, and recovery behavior. ProteinGym uses the official $217$-assay Average Spearman, VenusMutHub uses assay-macro averages, and VenusViroHub first averages within phenotype--backbone cells and then equally across its seven phenotype classes. Appendix~\ref{app:evaluation_details} specifies validation and uncertainty.

\subsection{A shared readout transfers across architectures and benchmarks}
\label{sec:transfer}

\paragraph{One operator transfers across model families and data regimes.}
Using this shared $59$-configuration panel, \textsc{REM2} raises mean ProteinGym Average Spearman from $0.381$ to $0.452$, VenusMutHub assay-macro Spearman from $0.154$ to $0.213$, and VenusViroHub hierarchical Spearman from $0.177$ to $0.270$. Every configuration improves on every benchmark. Gains are largest for autoregressive models on ProteinGym and VenusMutHub and for masked models on VenusViroHub, while structure-aware and inverse-folding models also benefit (Appendix~\ref{app:family_results}). Family evidence therefore contributes across the tested architectural range.

\paragraph{The gain spans complementary notions of ranking quality.}
Across Spearman, NDCG, AUC, MCC, and top-variant recall, \textsc{REM2} improves all $295$ model--metric comparisons on ProteinGym and VenusViroHub and $293$ of $295$ on VenusMutHub (Figure~\ref{fig:scatter}). The common direction across rank correlation, discrimination, and top-candidate recovery shows that the readout reshapes the mutation landscape throughout its dynamic range.

\paragraph{Gains follow a functional hierarchy.}
On ProteinGym, all $59$ configurations improve in each of the five functional categories ($295/295$ configuration--category comparisons). Mean gains are largest for Stability ($0.102$), Activity ($0.082$), and Organismal Fitness ($0.072$), followed by Binding ($0.050$) and Expression ($0.048$). The ordering concentrates the largest corrections in phenotypes most directly governed by conservation and folding.

\paragraph{\textsc{VenusREM2} sets a new ProteinGym state of the art.}
The six-vocabulary \textsc{VenusREM2} ensemble reaches an official Average Spearman of $0.556$, exceeding its uncalibrated counterpart by $0.031$ and the strongest reported public result by $0.038$. It ranks first in all five functional categories, reaching $0.691$ on Stability and $0.494$ on Organismal Fitness (Table~\ref{tab:leaderboard}).

\subsection{Scale separation predicts when family retrieval matters}
\label{sec:results_scales}

\begin{wrapfigure}[16]{R}{0.46\textwidth}
\vspace{-1.1em}
\centering
\includegraphics[width=\linewidth]{figures/fig_separation_gain.pdf}
\vspace{-1.2em}
\caption{Separation predicts MSA retrieval gain across $148$ ProteinGym proteins; red points are quartile means with 95\% bootstrap intervals.}
\label{fig:separation_gain}
\vspace{-1.4em}
\end{wrapfigure}

\paragraph{Separation identifies proteins that benefit from family evidence.}
Across $148$ proteins, the fraction of separated sites correlates with MSA retrieval gain at Spearman $\rho=0.338$, with bootstrap interval $0.179$--$0.481$ (Figure~\ref{fig:separation_gain}). Mean gain rises from $0.0037$ in the lowest-separation quartile to $0.0619$ in the highest, and the fraction of proteins that improve rises from $54.1\%$ to $83.8\%$. Retrieval is therefore most effective when homologs alter the native amino-acid ordering rather than merely reinforce it.

\paragraph{Direction carries information beyond alignment quantity.}
The Q4--Q1 gain difference is $0.058$, with bootstrap interval $0.033$--$0.086$. After controlling for sequence length, MSA depth, occupancy, and Raw performance, separation retains a partial rank correlation of $\rho=0.261$. Family history is thus most valuable where it contributes a new local direction, rather than a deeper sample of the ordering already encoded by the backbone.

\subsection{Retrieval supplies direction; calibration regulates confidence}
\label{sec:components}

\begin{table}[!t]
\setlength{\abovecaptionskip}{4pt}
\setlength{\belowcaptionskip}{2pt}
\caption{\textsc{ProSST}-2048 component ablation on ProteinGym validation/full, VenusMutHub, and VenusViroHub. Ent. Weight., Coh. Gating, and RSA denote entropy weighting, coherence gating, and relative solvent accessibility. w/o Ent. Weight. fixes $\alpha=0.8$; w/o Coh. Gating fixes $\kappa=1$, whereas w/ Coh. Gating uses $\kappa(R)$.}
\label{tab:ablation}
\centering
\scriptsize
\setlength{\tabcolsep}{1.8pt}
\resizebox{\textwidth}{!}{%
\begin{tabular}{@{}llcccccccc@{}}
\toprule
\multirow[c]{2}{*}{Component} & \multirow[c]{2}{*}{Setting} & \multicolumn{2}{c}{PG Validation (21)} & \multicolumn{2}{c}{PG Full (217)} & \multicolumn{2}{c}{VMH Full (905)} & \multicolumn{2}{c}{VVH (89)} \\
\cmidrule(lr){3-4}\cmidrule(lr){5-6}\cmidrule(lr){7-8}\cmidrule(l){9-10}
& & Score & $\Delta_{\mathrm{Raw}}$ & Score & $\Delta_{\mathrm{Raw}}$ & Score & $\Delta_{\mathrm{Raw}}$ & Score & $\Delta_{\mathrm{Raw}}$ \\
\midrule
Raw & Native field & 0.5573 & --- & 0.5069 & --- & 0.2126 & --- & 0.2944 & --- \\
\midrule
\multirow[c]{2}{*}{MSA retrieval} & w/o Ent. Weight. & 0.5742 & $+1.7\%$ & 0.5205 & $+1.4\%$ & \textbf{0.2341} & $\mathbf{+2.2\%}$ & 0.3114 & $+1.7\%$ \\
& w/ Ent. Weight. & \textbf{0.5774} & $\mathbf{+2.0\%}$ & \textbf{0.5246} & $\mathbf{+1.8\%}$ & 0.2336 & $+2.1\%$ & \textbf{0.3117} & $\mathbf{+1.7\%}$ \\
\midrule
\multirow[c]{2}{*}{Adaptive calibration} & w/o Coh. Gating & 0.5641 & $+0.7\%$ & 0.5219 & $+1.5\%$ & 0.2326 & $+2.0\%$ & 0.3062 & $+1.2\%$ \\
& w/ Coh. Gating & \textbf{0.5788} & $\mathbf{+2.2\%}$ & \textbf{0.5284} & $\mathbf{+2.2\%}$ & \textbf{0.2369} & $\mathbf{+2.4\%}$ & \textbf{0.3155} & $\mathbf{+2.1\%}$ \\
\midrule
\multirow[c]{2}{*}{Shrinkage} & w/ RSA, w/o pLDDT & 0.5826 & $+2.5\%$ & \textbf{0.5329} & $\mathbf{+2.6\%}$ & \textbf{0.2499} & $\mathbf{+3.7\%}$ & \textbf{0.3292} & $\mathbf{+3.5\%}$ \\
& w/o RSA, w/ pLDDT & \textbf{0.5833} & $\mathbf{+2.6\%}$ & \textbf{0.5329} & $\mathbf{+2.6\%}$ & 0.2352 & $+2.3\%$ & 0.3196 & $+2.5\%$ \\
\midrule
\textsc{REM2} & Full pipeline & \textbf{0.5849} & $\mathbf{+2.8\%}$ & \textbf{0.5343} & $\mathbf{+2.7\%}$ & 0.2491 & $+3.7\%$ & \textbf{0.3302} & $\mathbf{+3.6\%}$ \\
\bottomrule
\end{tabular}%
}
\end{table}

\paragraph{Retrieval contributes most of the gain; later stages regulate it.}
On \textsc{ProSST}-2048, retrieval supplies $0.0201$ of the $0.0276$ validation gain; calibration and reliability add $0.0075$, raising Spearman from $0.5573$ to $0.5849$. Full ProteinGym and VenusViroHub improve at every stage, whereas VenusMutHub peaks at $0.2499$ with RSA-only shrinkage and the full pipeline remains within $0.001$ at $0.2491$ (Table~\ref{tab:ablation}). Retrieval therefore determines the dominant direction, while subsequent stages provide smaller, benchmark-dependent regulation. A fixed-share sweep forms a broad validation plateau over $\alpha\in[0.6,0.9]$, with the entropy-adaptive coefficient at or above the best fixed share (Appendix~\ref{app:alpha_sweep}).

\paragraph{The hierarchy recurs across architectures.}
Across the shared panel, retrieval contributes approximately $54\%$, $61\%$, and $44\%$ of the total gains on ProteinGym, VenusMutHub, and VenusViroHub. Full \textsc{REM2} raises mean Spearman by $0.071$ ($0.381$ to $0.452$), $0.059$ ($0.154$ to $0.213$), and $0.093$ ($0.177$ to $0.270$), respectively. Recurrence across four backbone families and three data regimes supports the same division of labor (Appendix~\ref{app:readout_runtime}).

\paragraph{Coherence selectively licenses global calibration.}
Relative to retrieval alone, gated calibration improves $59/59$, $56/59$, and $59/59$ configurations on ProteinGym, VenusMutHub, and VenusViroHub; relative to ungated calibration, the corresponding counts are $59/59$, $23/59$, and $55/59$. The leave-one-site-out gate agrees with the original estimate (Spearman $\rho=0.985$; Pearson $r=0.988$), and gains concentrate in high-coherence quartiles. RSA and pLDDT then attenuate corrections at exposed or structurally uncertain sites (Appendix~\ref{app:bias_details}; Figure~\ref{fig:coherence_mechanism}).

\begin{figure}[!t]
\centering
\includegraphics[width=\textwidth]{figures/fig_viral_parkin.pdf}
\caption{\textbf{(a)} Measured escape versus frozen-backbone constraint on XBB.1.5 RBD and ZIKV E. \textbf{(b)} VenusViroHub phenotype-wise \textsc{REM2}$-$Raw gains; red marks unresolved regimes. \textbf{(c)} Parkin domain context: RING1, REP, and case sites. \textbf{(d)} Domain-wise Spearman. \textbf{(e)} Predicted percentiles move toward the measurements for N232A and R396H.}
\label{fig:viral_parkin}
\end{figure}

\subsection{Immune escape separates intrinsic fitness from host context}
\label{sec:immune_escape}

\paragraph{VenusViroHub gains are phenotype selective.}
Although \textsc{REM2} raises overall VenusViroHub hierarchical Spearman from $0.177$ to $0.270$, the gain is largest for expression ($+0.164$) and activity ($+0.145$), negligible for immune escape ($+0.003$; $0.025$ to $0.028$), and absent for stability ($0.070$ for both Raw and \textsc{REM2}; Figure~\ref{fig:viral_parkin}b). Across $21$ immune-escape assays and $12$ viral backbones, the best mean result is only $0.054$. Family calibration therefore improves viral mutation ranking broadly while leaving immune-escape ordering unresolved.

\paragraph{Immune selection can oppose evolutionary fitness.}
Among the $21$ escape assays, $11$ show positive signal (up to $+0.23$), $6$ remain near zero, and $4$ correlate negatively (down to $-0.30$). The negative regime localizes to known serum and monoclonal-antibody epitopes: on the SARS-CoV-2 RBD and ZIKV E, measured escape concentrates on exposed receptor-binding and antibody-contact ridges, while frozen backbones score the same sites as deleterious (Figure~\ref{fig:viral_parkin}a). Per-site escape and model tolerance consequently anti-correlate, reaching $-0.19$ for the strongest structure-aware backbones. Evolutionary fields capture intrinsic viability; immune-aware ranking additionally requires antibody, serum, or host context.

\subsection{Family evidence reorganizes mutation landscapes along protein architecture}
\label{sec:parkin_case}

\paragraph{Gains resolve from assay level to protein anatomy.}
On Parkin, Raw, MSA retrieval, and \textsc{REM2} reach Spearman correlations of $0.503$, $0.585$, and $0.645$ across $8{,}756$ substitutions at $462$ sites; abundance-class AUC rises from $0.769$ to $0.815$ and $0.857$ (Figure~\ref{fig:viral_parkin}c--e).

\paragraph{The corrections follow domain organization and act in both directions.}
\textsc{REM2} raises domain-resolved Spearman from $0.425$ to $0.627$ in RING0, $0.427$ to $0.529$ in RING2, and $-0.083$ to $0.221$ in REP/tether; the flexible Ubl--RING0 linker remains weakly conserved (Figure~\ref{fig:viral_parkin}c,d). This asymmetry recurs across the benchmark: within-region gains concentrate in ordered regions, while low-confidence regions mark a scope boundary of the readout (Appendix~\ref{app:parkin_details}). N232A moves from the Raw $88.8$th predicted percentile toward its measured $0.4$th percentile, reaching $37.3$ after \textsc{REM2}; R396H moves oppositely from $39.7$ to $78.0$, matching the measured $77.6$th percentile (Figure~\ref{fig:viral_parkin}e). These bidirectional shifts follow local evolutionary context.

\subsection{Scoring interfaces expose architecture-specific accuracy--efficiency trade-offs}
\label{sec:efficient_inference}

\begin{wrapfigure}{R}{0.44\textwidth}
\vspace{-1.0em}
\centering
\includegraphics[width=\linewidth]{figures/fig_runtime_main.pdf}
\vspace{-0.8em}
\caption{Runtime across $2.47$ million ProteinGym variants: $20.7\times$ faster mutation lookup and $62.3\times$ faster MSA construction.}
\label{fig:runtime}
\vspace{-1.0em}
\end{wrapfigure}

\paragraph{\textsc{REM2} gains persist across encoder conditioning conventions.}
Across eight paired encoder checkpoints, both masked and wild-type fields improve on ProteinGym, VenusMutHub, and VenusViroHub. All six dataset--interface aggregates increase, with gains from $0.045$ to $0.168$, showing that \textsc{REM2} is not tied to whether the native residue is visible during inference (Appendix~\ref{app:readout_runtime}).

\paragraph{The preferred scoring interface depends on model architecture.}
On ProteinGym, full-sequence and native-field scoring are nearly equivalent for \textsc{ESM-IF1} ($0.422$ versus $0.421$), whereas full-sequence likelihood leads by $0.040$--$0.053$ across five \textsc{ProGen2} scales. \textsc{ProteinMPNN} reverses this pattern: removing random decoding-order variation raises mean assay Spearman from $0.271$--$0.290$ to $0.376$--$0.401$ across eight checkpoints. These contrasts establish an architecture-dependent optimum: native fields provide the common \textsc{REM2} substrate, while full-sequence likelihood can retain useful autoregressive context (Appendix~\ref{app:readout_runtime}). The sign reversal separates two sources of autoregressive context: downstream likelihood carries useful sequence coupling in \textsc{ProGen2}, whereas randomized decoding orders inject variance in \textsc{ProteinMPNN}.

\paragraph{A reusable field removes mutation-wise computation.}
For \textsc{ProteinMPNN}, native teacher forcing couples the accuracy gain above with a $219$--$601\times$ reduction in wall-clock cost. After a protein field is constructed, indexed lookup scores $2.47$ million ProteinGym variants $20.7\times$ faster, while vectorized MSA-field construction is $62.3\times$ faster (Figure~\ref{fig:runtime}). A simpler shared scoring field, GPU-resident execution, batched tokenization, vectorized construction, and indexed lookup remove repeated mutation-wise work.

\paragraph{Teacher forcing sharpens main effects without recovering epistasis.}
On the complete Olson GB1 mutation quartets, native teacher forcing raises single-mutant Spearman from $0.105$ to $0.306$ and double-mutant Spearman from $0.165$ to $0.284$. Its additive field contains no pairwise interaction by construction; allowing background-dependent conditional rescoring yields an epistasis correlation of only $0.012$, close to $0.005$ under the official score. The \textsc{ProteinMPNN} improvement therefore reflects a less noisy estimate of first-order substitution effects, not latent recovery of higher-order coupling. Epistasis remains a separate modeling target rather than a consequence of score-factorization choice (Appendix~\ref{app:readout_runtime}).

\FloatBarrier

\section{Discussion}
\label{sec:discussion}

\textsc{REM2} supports scale-separated mutation prediction: frozen backbones encode global amino-acid grammar, homologs supply local direction, and coherence and structure regulate calibration and uncertainty. Transfer across $59$ configurations and three benchmarks, together with the Parkin and VenusViroHub analyses, connects these gains to domains, residues, and phenotype context. Native teacher forcing exposes a reusable main-effect field but leaves GB1 epistasis near zero (Appendix~\ref{app:readout_runtime}), motivating post-training or preference learning for family- and host-conditioned effects.

\subsection*{AI Use Statement}
In this work, generative AI tools (OpenAI Codex) were used to refine the conceptual framing and hypotheses; to provide feedback on methodology and experimental design; to implement methods; to reformat data; to interpret and present experimental analyses; and to assist with translation. They were not used to generate synthetic datasets; theorem- and proof-writing tasks are not applicable to this work. Additionally, they were used to create and edit scientific figures, to edit code, to search and summarize literature, and to structure, draft, and language-edit the manuscript. The authors reviewed all AI-assisted work, verified reported results against the underlying evaluation artifacts and citations against their primary sources, and take full responsibility for the final text, claims, code, and figures.

\subsection*{Reproducibility Statement}
Sections~\ref{sec:method}--\ref{sec:ranking} specify the readout and scoring equations, while Section~\ref{sec:benchmarks} and Appendix~\ref{app:evaluation_details} document evaluation, aggregation, validation, and uncertainty estimation. Appendices~\ref{app:baselines}--\ref{app:replications} provide the baseline taxonomy, native-background diagnostics, full cross-family results, adaptive-component and gate ablations, interface and runtime audits, per-configuration metrics, the VenusViroHub benchmark specification, and the VenusMutHub replication. An anonymized implementation snapshot and the run configurations needed to reproduce the reported tables and figures are given in the Code Availability statement below.

\subsection*{Ethics Statement}
This study uses public protein sequence, structure, alignment, and deep-mutational-scanning benchmark data and involves no human participants or personal data. Mutation-effect scores are computational hypotheses that require experimental validation; reuse should follow the licenses and terms of the original datasets and pretrained models.

\subsection*{Code Availability}
The implementation, evaluation scripts, and experiment configurations are archived at:\par
\noindent\href{https://github.com/anonymous/VenusREM2}{\texttt{github.com/anonymous/VenusREM2}}.

\FloatBarrier

\setlength{\bibsep}{0pt plus 0.3ex}
\bibliography{0references}
\bibliographystyle{iclr2027_conference}

\clearpage
\appendix

% Appendix-only float tuning: favor compact top placement without shrinking content.
\setcounter{topnumber}{4}
\setcounter{bottomnumber}{2}
\setcounter{totalnumber}{6}
\renewcommand{\topfraction}{0.97}
\renewcommand{\bottomfraction}{0.90}
\renewcommand{\textfraction}{0.03}
\renewcommand{\floatpagefraction}{0.72}
\setlength{\textfloatsep}{8pt plus 2pt minus 2pt}
\setlength{\floatsep}{7pt plus 2pt minus 2pt}
\setlength{\intextsep}{8pt plus 2pt minus 2pt}

\section{Related Work}
\label{app:related_work}

\paragraph{Protein models.}
Protein mutation-effect prediction combines family-specific evolutionary models with pretrained protein models. \textsc{EVmutation}, \textsc{GEMME}, and \textsc{EVE} infer homolog-family constraints, while \textsc{RSALOR} and \textsc{ESCOTT} add structural context \citep{hopf2017evmutation,laine2019gemme,frazer2021eve,tsishyn2025rsalor,tekpinar2025prescott}. ProteinGym standardizes zero-shot evaluation \citep{notin2024proteingym} across masked and autoregressive sequence models \citep{meier2021esm1v,rives2021esm1b,lin2023esm2,madani2023progen,nijkamp2023progen2,hesslow2022rita,carp}, structure-aware encoders \citep{su2023saprot,li2024prosst,hayes2025esm3,zhang2024s3f,tan2025protssn}, and inverse-folding models \citep{hsu2022esm-if1,dauparas2022proteinmpnn,mifst}. These families supply complementary sequence, evolutionary, and geometric priors.

\paragraph{Post-hoc calibration.}
Post-hoc methods add evidence without redesigning the backbone. Retrieval and family modeling underlie \textsc{MSA Transformer}, \textsc{Tranception}/\textsc{TranceptEVE}, \textsc{PoET}, \textsc{Protriever}, \textsc{S3F-MSA}/\textsc{S2F-MSA}, and \textsc{VenusREM} \citep{rao2021msa,notin2022tranception,notin2022trancepteve,truong2023poet,notin2025protriever,zhang2024s3f,tan2025venusrem}. Supervised alternatives fine-tune protein-specific predictors, post-train pretrained representations, or learn multimodal ensemble weights \citep{zhou2024protlgn,ouyang2024predicting,tan2026rank}; homolog fine-tuning further links accuracy to wild-type likelihood and family coverage \citep{gordon2025preference}. \textsc{REM2} instead combines family constraints, frozen scores, and structural reliability in one training-free, architecture-agnostic readout.

\section{Evaluation, validation, and provenance details}
\label{app:evaluation_details}

\subsection{Metric definitions and benchmark aggregation}
For an assay with $n$ variants retained after removing non-finite pairs, let $y_i$ be the experimental effect and $s_i$ the model score. Let $u_i=\operatorname{rank}(y_i)$ and $v_i=\operatorname{rank}(s_i)$ use average ranks for ties. Assay-level Spearman correlation is
\begin{equation}
\rho_{\mathrm{S}}(y,s)=
\frac{\sum_{i=1}^{n}(u_i-\bar u)(v_i-\bar v)}
{\sqrt{\sum_{i=1}^{n}(u_i-\bar u)^2}\sqrt{\sum_{i=1}^{n}(v_i-\bar v)^2}}.
\label{eq:metric_spearman}
\end{equation}

For NDCG@10\%, define linear min--max gains $g_i=(y_i-\min_j y_j)/(\max_j y_j-\min_j y_j)$, $k=\lfloor0.1n\rfloor$, and descending ranks $r_i^{(s)}$ and $r_i^{(y)}$ under $s$ and $y$. The predicted and ideal discounted gains are
\begin{equation}
\begin{aligned}
\mathrm{DCG}_{10}
&=\sum_{i:r_i^{(s)}\le k}\frac{g_i}{\log_2(r_i^{(s)}+1)}, &
\mathrm{IDCG}_{10}
&=\sum_{i:r_i^{(y)}\le k}\frac{g_i}{\log_2(r_i^{(y)}+1)},\\
\mathrm{NDCG}_{10}
&=\frac{\mathrm{DCG}_{10}}{\mathrm{IDCG}_{10}}. & &
\end{aligned}
\label{eq:metric_ndcg}
\end{equation}
Let $Q_{0.9}(\cdot)$ denote the empirical 90th percentile and define $T_y=\{i:y_i\ge Q_{0.9}(y)\}$ and $T_s=\{i:s_i\ge Q_{0.9}(s)\}$. Top-variant recall@10\% is
\begin{equation}
\mathrm{Recall}_{10}=\frac{|T_y\cap T_s|}{|T_y|}.
\label{eq:metric_toprecall}
\end{equation}

For binary discrimination, $z_i=\mathbf{1}\{y_i\ge\operatorname{median}(y)\}$ defines the experimental positive class, with $n_+=\sum_i z_i$ and $n_-=n-n_+$. AUC uses the continuous model scores and average treatment of ties,
\begin{equation}
\mathrm{AUC}=\frac{1}{n_+n_-}\sum_{i:z_i=1}\sum_{j:z_j=0}
\left[\mathbf{1}\{s_i>s_j\}+\tfrac12\mathbf{1}\{s_i=s_j\}\right].
\label{eq:metric_auc}
\end{equation}
For MCC, predicted labels are $\hat z_i=\mathbf{1}\{s_i\ge\operatorname{median}(s)\}$, giving
\begin{equation}
\mathrm{MCC}=\frac{TP\,TN-FP\,FN}
{\sqrt{(TP+FP)(TP+FN)(TN+FP)(TN+FN)}}.
\label{eq:metric_mcc}
\end{equation}
Metrics are undefined for constant rankings, empty ideal gain, or degenerate binary classes and are omitted from subsequent means. All metrics are computed per assay before benchmark aggregation. VenusMutHub reports the assay macro-average. For VenusViroHub, if $m_{pba}$ is the assay metric for phenotype $p$, backbone $b$, and assay $a$, then
\begin{equation}
M_{\mathrm{VenusViroHub}}=\frac{1}{|\mathcal P|}\sum_{p\in\mathcal P}
\frac{1}{|\mathcal B_p|}\sum_{b\in\mathcal B_p}
\frac{1}{|\mathcal A_{pb}|}\sum_{a\in\mathcal A_{pb}}m_{pba},
\label{eq:virogym_aggregation}
\end{equation}
so densely assayed backbones do not dominate. ProteinGym leaderboard Spearman follows the official 217-assay evaluator; extended metrics use the same per-assay definitions above.

Official Average Spearman, VenusMutHub macro-averaging, and VenusViroHub hierarchical aggregation follow Section~\ref{sec:benchmarks}. Uncertainty uses 20{,}000 paired bootstrap draws. For ProteinGym we resample 186 unique UniProt IDs and include every assay and selection type of a sampled protein; for VenusMutHub we resample 527 unique wild-type (WT) sequence clusters and include every assay in a sampled cluster; for VenusViroHub we resample 52 phenotype--backbone cells within each of the seven phenotype strata and reapply the hierarchy. The same draws are applied to Raw and \textsc{REM2}. The 59 configurations are held fixed, so the intervals quantify uncertainty over benchmark units, not over a model population. In Figure~\ref{fig:motivation}b, each well-covered site contributes Pearson correlations across the 20 amino-acid entries. Duplicate assays are collapsed by WT sequence, yielding $27{,}763$, $148{,}983$, and $29{,}416$ complete profiles from $148$, $527$, and $52$ representative sequences in ProteinGym, VenusMutHub, and VenusViroHub, respectively. The plot uses $7{,}500$ profiles per benchmark sampled with seed 20260904; all summaries use the complete sets. The cutoff $r=0.3$ is descriptive; the claim rests on the continuous distributions and their medians.

In Figure~\ref{fig:separation_gain}, the protein-level separation index is the fraction of well-covered sites---those with complete 20-amino-acid MSA and native contrast vectors---whose MSA preference is not jointly aligned with log UniProt abundance and the protein-mean native profile at $r\geq0.3$. All 217 \textsc{ESM-2} wild-type-interface ProteinGym assays have paired Raw--MSA retrieval scores; 177 contain at least 20 complete single-site amino-acid profiles, spanning 148 unique proteins. Assays that share a protein are aggregated before the protein-level analysis: separation and MSA gain are averaged within protein, and the site-level map uses, for each protein, only the assay with the broadest complete profile. MSA gain is the protein-mean paired assay Spearman difference. 10{,}000 protein-level bootstrap draws give the intervals in Figure~\ref{fig:separation_gain}; partial rank correlation controls for sequence length, $\log_{10}$ MSA depth, MSA occupancy, and Raw Spearman. The cross-benchmark site-level replication uses the same \textsc{ESM-2} 650M wild-type interface and is independent of assay labels.

All ProteinGym results come from one versioned evaluation pipeline. Of 60 configurations, 59 have complete paired Raw--\textsc{REM2} coverage on 217 assays; the \textsc{VenusREM} configuration was not rescored with the \textsc{REM2} readout and is excluded from paired comparisons. A2M files are aligned FASTA with match states; lowercase letters are converted to uppercase, and gaps are mapped to padding. Retained rows are counted equally, without sequence reweighting or pseudocounts. The general operator is the entropy-weighted protein-wise map in Equation~\ref{eq:native_alpha}, applied independently to every backbone in the leaderboard and family tables. Component operators are selected on the 21-assay validation split---fold~8 of a ten-fold partition formed after shuffling the sorted assay identifiers with random seed~42---and then frozen; the primary ProteinGym result uses official Average Spearman over all 217 assays.

For leave-one-site-out coherence analysis, site $i$ uses $b_a^{(-i)}=\operatorname{logsumexp}_{j\ne i}R_{j,a}-\log(L-1)$, and the leave-one-site-out gate is $\kappa_{\mathrm{LOSO}}=[c_{\mathrm{LOSO}}]_+$, where $c_{\mathrm{LOSO}}=L^{-1}\sum_i\operatorname{corr}_a(R_i,b^{(-i)})$. The paired gain is the assay-level Spearman difference between adaptive calibration (weight $1-\alpha$) and MSA retrieval from the same \textsc{ProSST-2048} field. Across 217 assays and 187 unique wild-type sequences, Figure~\ref{fig:coherence_mechanism}b reports assay-bootstrap intervals within coherence quartiles.

\subsection{Dataset and input distributions}

\paragraph{Sequence and alignment distributions.}
ProteinGym contains 217 assays spanning 186 unique UniProt IDs and 187 unique wild-type (WT) sequences, VenusMutHub contains 905 assays across 527 unique WT clusters, and VenusViroHub contributes 89 assays across 52 unique WT sequences after duplicate removal. Figure~\ref{fig:dataset_stats}a--c and Table~\ref{tab:dataset_stats} compare their inputs. VenusViroHub sequences are longer on average (559 residues versus 397 and 288) but have shallower MSAs (1{,}217 sequences versus 8{,}650 and 9{,}791); mean MSA occupancy is 0.62, compared with 0.67 for ProteinGym and 0.69 for VenusMutHub.

\begin{figure}[!ht]
\centering
\includegraphics[width=\textwidth]{figures/fig_dataset_statistics.pdf}
\caption{Input distributions for ProteinGym (217 assays), VenusMutHub (905), and VenusViroHub (89). \textbf{(a)} Sequence length. \textbf{(b)} MSA depth. \textbf{(c)} MSA occupancy. \textbf{(d)} Residue pLDDT. \textbf{(e)} Residue RSA. Structure panels pool unique WT sequences; all selected structures fully cover the query.}
\label{fig:dataset_stats}
\end{figure}

\input{table_dataset_statistics.tex}

\paragraph{Structure distributions.}
pLDDT is read from C$\alpha$ B-factors and RSA from Shrake--Rupley accessibility normalized by residue-specific maximum ASA. Selected AlphaFold2 models cover every query residue in all three benchmarks. VenusViroHub has the lowest assay-level mean pLDDT (83.9), with 41.1\% of residues at pLDDT $\geq90$ and 9.7\% below 50, while its mean RSA (0.33) lies between ProteinGym (0.37) and VenusMutHub (0.32). Figure~\ref{fig:dataset_stats}d--e and Tables~\ref{tab:plddt_ratios}--\ref{tab:rsa_ratios} report the full distributions.

\input{table_dataset_structure_ratios.tex}

\subsection{VenusViroHub data processing, provenance, and dataset tables (TODO)}
\label{app:virogym_data_processing}

\noindent\textbf{TODO.} Add VenusViroHub source datasets and references, the data-cleaning and harmonization procedures, and dataset-level provenance and cleaning tables.

\subsection{VenusViroHub benchmark}
\label{app:virogym}

\paragraph{Scope and curation.}
VenusViroHub contains $90$ catalog entries collected from viral deep-mutational scans; one exact duplicate is excluded, leaving $89$ assays and $646{,}756$ scored substitutions. The benchmark spans $19$ viruses, $31$ backbone groups, and seven phenotype classes: immune escape, cell entry, binding, expression, fitness, stability, and activity. Its assays have zero overlap with ProteinGym.

\paragraph{Hierarchical evaluation.}
To prevent densely assayed viral backbones from dominating, each metric is first averaged within phenotype--backbone cells, then across backbones within a phenotype, and finally equally across the seven phenotypes. The $59$ Raw--\textsc{REM2} pairs use the same frozen checkpoints and five metric definitions as the other benchmarks.

\paragraph{Cross-model transfer.}
Under this hierarchy, mean Spearman rises from $0.177$ to $0.270$ ($+0.093$), with positive changes for all $59$ configurations. The largest gains occur for masked sequence models ($+0.140$; $0.106$ to $0.246$) and autoregressive models ($+0.107$; $0.121$ to $0.228$). Structure-aware encoders improve from $0.240$ to $0.301$ ($+0.061$), and inverse-folding models from $0.288$ to $0.325$ ($+0.037$). Thus, family evidence contributes most strongly to sequence-only readouts while remaining complementary to geometric conditioning (Table~\ref{tab:virogym_family}).

\paragraph{Metric and component consistency.}
Averaged over the same $59$ configurations, NDCG increases from $0.704$ to $0.753$, AUC from $0.601$ to $0.652$, MCC from $0.150$ to $0.225$, and top-variant recall from $0.142$ to $0.170$. Every configuration improves on every metric, yielding $295$ of $295$ positive model--metric comparisons (Tables~\ref{tab:viro_metrics_family_1}--\ref{tab:viro_metrics_family_4}). For \textsc{ProSST}-2048, entropy-weighted retrieval raises Spearman from $0.2944$ to $0.3117$, and coherence-gated calibration reaches $0.3155$. RSA-only and pLDDT-only shrinkage reach $0.3292$ and $0.3196$, respectively, while the full readout reaches $0.3302$, identifying exposure as the larger reliability correction on VenusViroHub (Table~\ref{tab:ablation}).

\FloatBarrier
\section{Baseline models and implementations}
\label{app:baselines}

Table~\ref{tab:baselines} lists the evaluated baselines, their readout interfaces, and official implementations. All backbones are frozen and evaluated with their native readouts (Equation~\ref{eq:readouts}); Appendix~\ref{app:bias_details} quantifies the shared native-background direction and interface-specific checkpoint deviations induced by these readouts. Configurations from the same family (e.g., the \textsc{ESM-2} size series) count as distinct model--readout configurations. The \textsc{ProSST} series varies the structure-vocabulary size $K$. \textsc{VenusREM2} scores each of the six $K$-variants with \textsc{REM2} at that variant's own entropy-weighted $\alpha$ from Equation~\ref{eq:native_alpha}, then $z$-scores each $K$'s per-assay variant-score vector and averages the six standardized vectors. \textsc{S3F} contributes two readout configurations, wild-type and masked marginals, from one checkpoint but with distinct native fields. \textsc{SaProt} is evaluated with PDB- and AlphaFold2-trained weights and with both interfaces. Protein structures are AlphaFold2 predictions \citep{jumper2021alphafold2} and query-matched alignments are built with ColabFold \citep{mirdita2022colabfold}.

\begin{table}[!ht]
\caption{Evaluated baselines, native readouts, and implementations. Readouts follow Equation~\ref{eq:readouts}; mask $=$ masked marginal, wt $=$ wild-type field, and tf $=$ teacher forcing.}
\label{tab:baselines}
\begin{center}
\scriptsize
\setlength{\tabcolsep}{3pt}
\resizebox{\textwidth}{!}{%
\begin{tabular}{lllll}
\toprule
Model & Family & Evaluated variants & Readout & Source \\
\midrule
\textsc{ESM-1b} & Masked sequence LM & 650M & mask, wt & \href{https://github.com/facebookresearch/esm}{Official repository} \\
\textsc{ESM-1v} & Masked sequence LM & 650M & mask, wt & \href{https://github.com/facebookresearch/esm}{Official repository} \\
\textsc{ESM-2} & Masked sequence LM & 8M/35M/150M/650M/3B & mask, wt & \href{https://github.com/facebookresearch/esm}{Official repository} \\
\textsc{ESMC} & Masked sequence LM & 300M, 600M & mask & \href{https://github.com/evolutionaryscale/esm}{Official repository} \\
\textsc{CARP} & Masked sequence LM & 640M & mask & \href{https://github.com/microsoft/protein-sequence-models}{Official repository} \\
\textsc{S3F} & Structure-aware & S3F & wt, mask & \href{https://github.com/DeepGraphLearning/S3F}{Official repository} \\
\textsc{ProGen2} & Autoregressive LM & B/S/M/L/XL & tf & \href{https://github.com/salesforce/progen}{Official repository} \\
\textsc{ProGen3} & Autoregressive LM & 112M/219M/339M/762M/1B/3B & tf & \href{https://github.com/Profluent-AI/progen3}{Official repository} \\
\textsc{RITA} & Autoregressive LM & S/M/L/XL & tf & \href{https://github.com/lightonai/rita}{Official repository} \\
\textsc{ESM3} & Structure-aware & open-weight (1.4B) & mask & \href{https://github.com/evolutionaryscale/esm}{Official repository} \\
\textsc{ProSST} & Structure-aware & $K{=}20/128/512/1024/2048/4096$ & mask & \href{https://github.com/ai4protein/ProSST}{Official repository} \\
\textsc{ProtSSN} & Structure-aware & ensemble & tf & \href{https://github.com/ai4protein/ProtSSN}{Official repository} \\
\textsc{SaProt} & Structure-aware & 35M-AF2, 650M-PDB & mask, wt & \href{https://github.com/westlake-repl/SaProt}{Official repository} \\
\textsc{ESM-IF1} & Inverse folding & --- & tf & \href{https://github.com/facebookresearch/esm}{Official repository} \\
\textsc{ProteinMPNN} & Inverse folding & v\_48\_002/010/020/030; ProteinMPNN-Soluble $\times$4 & tf & \href{https://github.com/dauparas/ProteinMPNN}{Official repository} \\
\textsc{MIF-ST} & Inverse folding & --- & tf & \href{https://github.com/microsoft/protein-sequence-models}{Official repository} \\
\textsc{VenusREM} & Retrieval-enhanced & \textsc{ProSST} ($K{=}2048$) + MSA & mask & \href{https://github.com/ai4protein/VenusREM}{Official repository} \\
\textsc{RSALOR} & MSA + structure & conservation $\times$ RSA & --- & \href{https://github.com/3BioCompBio/RSALOR}{Official repository} \\
\bottomrule
\end{tabular}}%
\end{center}
\end{table}

The compact leaderboard in Table~\ref{tab:leaderboard} keeps one official row per model series. Table~\ref{tab:leaderboard_full} restores the remaining official entries used for that selection. Table~\ref{tab:leaderboard_full_taxon_msa} keeps those same rows and replaces the five biological-category columns with official taxon and MSA-depth bins: \textsc{VenusREM2} is strongest in every reported stratum, including virus ($0.540$) and the low-depth bin ($0.553$).

\input{table1_leaderboard_full.tex}

\input{table1_leaderboard_full_taxon_msa.tex}

\FloatBarrier
\section{Native-background diagnostics and checkpoint deviations}
\label{app:bias_details}

Figure~\ref{fig:bias} expands the consensus view in Figure~\ref{fig:teaser}: Figure~\ref{fig:bias}a overlays the 37 checkpoint profiles, Figure~\ref{fig:bias}b quantifies their pairwise agreement, Figure~\ref{fig:bias}c separates readout scales, and Figure~\ref{fig:bias}d shows that mean directional alignment $c$ spans $0.01$--$0.99$ across protein--configuration pairs. The profiles have a median pairwise Pearson correlation of $0.995$, and $0.993$ across families, and their consensus tracks log UniProtKB/Swiss-Prot amino-acid abundance with $r=0.99$. The spread in Figure~\ref{fig:bias}d motivates the adaptive gate. Figure~\ref{fig:coherence_mechanism}a shows that leave-one-site-out estimation preserves the gate (Spearman $\rho=0.985$; Pearson $r=0.988$), with mean $|\Delta\kappa|=0.018$ and median $|\Delta\kappa|=0.010$. Figure~\ref{fig:coherence_mechanism}b relates this independent estimate to calibration gain.

\textsc{ESM3} is the clearest checkpoint-level departure from the consensus, consistent with pretraining on a mixed corpus that includes text and other non-sequence modalities \citep{hayes2025esm3}. Differences across scoring interfaces of the same checkpoint are smaller, implicating pretraining mixture as the primary axis for future controlled studies of corpus, modality, tokenization, and readout.

\begin{figure}[!ht]
\begin{center}
\includegraphics[width=\textwidth]{figures/fig_bias.png}
\end{center}
\caption{Native-background diagnostics across 37 checkpoints. \textbf{(a)} Protein-level amino-acid profiles versus UniProtKB/Swiss-Prot composition. \textbf{(b)} Pairwise Pearson correlations among those profiles. \textbf{(c)} Readout-scale comparison across scoring interfaces. \textbf{(d)} Mean directional alignment $c=L^{-1}\sum_i\operatorname{corr}_a(R_i,b)$ over protein--configuration pairs.}
\label{fig:bias}
\end{figure}

The gate is therefore a protein-level property of $R$, not an assay-label fit. Figure~\ref{fig:coherence_mechanism}a shows that leaving site $i$ out of $b$ barely moves $\kappa$; Figure~\ref{fig:coherence_mechanism}b then places the calibration gain against that independent coherence estimate, with the upper quartiles carrying the mean lift.

\begin{figure}[!ht]
\centering
\includegraphics[width=\textwidth]{figures/fig_coherence_mechanism.pdf}
\caption{Leave-one-site-out coherence and calibration gain across 217 \textsc{ProSST-2048} assays. \textbf{(a)} The leave-one-site-out gate $\kappa_{\mathrm{LOSO}}$ preserves the full-protein gate $\kappa$. \textbf{(b)} Assay-level Spearman gain of adaptive calibration (weight $1-\alpha$) over MSA retrieval versus coherence; red points are quartile means with 95\% assay-bootstrap intervals.}
\label{fig:coherence_mechanism}
\end{figure}

\FloatBarrier
\section{Full cross-family results}
\label{app:family_results}

The 59-configuration ProteinGym comparison is uniform in sign and uneven in magnitude (Table~\ref{tab:proteingym}): autoregressive models gain the most, and structure-aware encoders still gain after geometry is already in the backbone. Table~\ref{tab:prosst_metrics} shows that the \textsc{ProSST} ensemble improves every reported metric, not only Spearman.

\begin{table}[!ht]
\caption{ProteinGym Average Spearman by backbone family across 59 paired configurations and 217 assays. $N$ is the number of configurations. Mean and median change are paired \textsc{REM2}--Raw differences; Wins / $N$ counts configurations with positive change.}
\label{tab:proteingym}
\begin{center}
\footnotesize
\setlength{\tabcolsep}{4pt}
\begin{tabular}{lcccccc}
\toprule
Family & $N$ & Raw & \textsc{REM2} & Mean change & Median change & Wins / $N$ \\
\midrule
Structure-aware & 17 & 0.456 & 0.495 & +0.039 & +0.031 & 17 / 17 \\
Masked sequence LM & 17 & 0.374 & 0.447 & +0.073 & +0.058 & 17 / 17 \\
Inverse folding & 10 & 0.382 & 0.447 & +0.065 & +0.070 & 10 / 10 \\
Autoregressive LM & 15 & 0.302 & 0.410 & +0.108 & +0.098 & 15 / 15 \\
\bottomrule
\end{tabular}
\end{center}
\end{table}

\begin{table}[!ht]
\caption{\textsc{ProSST}-ensemble ProteinGym metrics, reported as assay-macro mean $\pm$ standard deviation over 217 assays. These means are not official Average Spearman. $\Delta$ is \textsc{REM2}--Raw; the \textsc{REM2} column is the six-$K$ \textsc{VenusREM2} ensemble; top-recall denotes top-variant recall.}
\label{tab:prosst_metrics}
\centering
\small
\setlength{\tabcolsep}{9pt}
\begin{tabular}{lccc}
\toprule
Metric & Raw & \textsc{REM2} & $\Delta$ \\
\midrule
Spearman & 0.541 $\pm$ 0.180 & 0.575 $\pm$ 0.159 & +0.034 \\
NDCG & 0.794 $\pm$ 0.138 & 0.813 $\pm$ 0.130 & +0.019 \\
AUC & 0.796 $\pm$ 0.100 & 0.815 $\pm$ 0.091 & +0.019 \\
MCC & 0.422 $\pm$ 0.164 & 0.450 $\pm$ 0.151 & +0.028 \\
Top-recall & 0.261 $\pm$ 0.131 & 0.270 $\pm$ 0.131 & +0.008 \\
\bottomrule
\end{tabular}
\end{table}

The aggregate gain is not confined to one of the five ProteinGym biological categories. Across those categories and $59$ configurations, all $295$ configuration--category comparisons improve in Spearman. Table~\ref{tab:function_types} reports the corresponding family means; Stability shows the largest overall increase, gaining $0.102$ from $0.482$ to $0.584$.

\begin{table}[!ht]
\caption{Mean ProteinGym Average Spearman by biological category and backbone family across the 59 configurations. Counts in the header are the number of configurations in that family. Structure $=$ structure-aware; Masked LM $=$ masked sequence LM; AR LM $=$ autoregressive LM.}
\label{tab:function_types}
\begin{center}
\scriptsize
\setlength{\tabcolsep}{2.0pt}
\resizebox{\textwidth}{!}{%
\begin{tabular}{lcccccccccc}
\toprule
Category & \multicolumn{2}{c}{Structure (17)} & \multicolumn{2}{c}{Masked LM (17)} & \multicolumn{2}{c}{Inverse folding (10)} & \multicolumn{2}{c}{AR LM (15)} & \multicolumn{2}{c}{All (59)} \\
\cmidrule(lr){2-3}\cmidrule(lr){4-5}\cmidrule(lr){6-7}\cmidrule(lr){8-9}\cmidrule(lr){10-11}
& Raw & \textsc{REM2} & Raw & \textsc{REM2} & Raw & \textsc{REM2} & Raw & \textsc{REM2} & Raw & \textsc{REM2} \\
\midrule
Activity & 0.436 & 0.490 & 0.387 & 0.461 & 0.318 & 0.414 & 0.321 & 0.434 & 0.373 & 0.454 \\
Binding & 0.395 & 0.421 & 0.300 & 0.360 & 0.327 & 0.384 & 0.251 & 0.313 & 0.319 & 0.370 \\
Expression & 0.472 & 0.498 & 0.385 & 0.437 & 0.382 & 0.429 & 0.334 & 0.400 & 0.396 & 0.444 \\
Organismal Fitness & 0.375 & 0.432 & 0.315 & 0.404 & 0.294 & 0.373 & 0.334 & 0.397 & 0.334 & 0.405 \\
Stability & 0.604 & 0.637 & 0.483 & 0.571 & 0.588 & 0.635 & 0.272 & 0.507 & 0.482 & 0.584 \\
\bottomrule
\end{tabular}}
\end{center}
\end{table}

VenusMutHub repeats the same family order (Table~\ref{tab:venusmuthub}): every family improves, and autoregressive models again gain the most. Table~\ref{tab:bootstrap} summarizes the paired benchmark units: UniProt IDs for ProteinGym, wild-type sequence clusters for VenusMutHub, and phenotype--backbone cells for VenusViroHub.

\begin{table}[!ht]
\caption{VenusMutHub macro-averaged assay-level Spearman by backbone family across 59 paired configurations and 905 assays. Column definitions match Table~\ref{tab:proteingym}.}
\label{tab:venusmuthub}
\begin{center}
\footnotesize
\setlength{\tabcolsep}{4pt}
\begin{tabular}{lcccccc}
\toprule
Family & $N$ & Raw & \textsc{REM2} & Mean change & Median change & Wins / $N$ \\
\midrule
Structure-aware & 17 & 0.176 & 0.228 & +0.051 & +0.037 & 17 / 17 \\
Masked sequence LM & 17 & 0.145 & 0.201 & +0.057 & +0.045 & 17 / 17 \\
Inverse folding & 10 & 0.210 & 0.256 & +0.046 & +0.049 & 10 / 10 \\
Autoregressive LM & 15 & 0.102 & 0.182 & +0.080 & +0.070 & 15 / 15 \\
\bottomrule
\end{tabular}
\end{center}
\end{table}

VenusViroHub extends the comparison to a virus-focused distribution with no ProteinGym assay overlap (Table~\ref{tab:virogym_family}). The largest mean lift occurs for masked sequence models, while every configuration improves in every family.

\begin{table}[!ht]
\caption{VenusViroHub hierarchically aggregated Spearman by backbone family across 59 paired configurations and 89 assays. Column definitions match Table~\ref{tab:proteingym}.}
\label{tab:virogym_family}
\begin{center}
\footnotesize
\setlength{\tabcolsep}{4pt}
\begin{tabular}{lcccccc}
\toprule
Family & $N$ & Raw & \textsc{REM2} & Mean change & Median change & Wins / $N$ \\
\midrule
Structure-aware & 16 & 0.240 & 0.301 & +0.061 & +0.065 & 16 / 16 \\
Masked sequence LM & 18 & 0.106 & 0.246 & +0.140 & +0.140 & 18 / 18 \\
Inverse folding & 10 & 0.288 & 0.325 & +0.037 & +0.037 & 10 / 10 \\
Autoregressive LM & 15 & 0.121 & 0.228 & +0.107 & +0.108 & 15 / 15 \\
\bottomrule
\end{tabular}
\end{center}
\end{table}


\begin{table}[!ht]
\caption{Mean paired \textsc{REM2}--Raw changes across 59 configurations. $\Delta$ follows each benchmark's primary aggregation: Average Spearman for ProteinGym, assay-macro Spearman for VenusMutHub, and hierarchical Spearman for VenusViroHub. Units follow the paired bootstrap in Appendix~\ref{app:evaluation_details}; Wins / $N$ counts positive changes.}
\label{tab:bootstrap}
\begin{center}
\small
\setlength{\tabcolsep}{4pt}
\begin{tabular}{llccc}
\toprule
Benchmark & Resampling unit & Units & Mean $\Delta$ & Wins / $N$ \\
\midrule
ProteinGym & UniProt ID & 186 & +0.071 & 59 / 59 \\
VenusMutHub & wild-type sequence & 527 & +0.059 & 59 / 59 \\
VenusViroHub & phenotype--backbone cell & 52 & +0.093 & 59 / 59 \\
\bottomrule
\end{tabular}
\end{center}
\end{table}

\FloatBarrier
\section{Parkin abundance landscape}
\label{app:parkin_details}

Parkin is a 465-residue RING-between-RING E3 ubiquitin ligase whose abundance landscape measures how missense variation preserves cellular proteostasis \citep{clausen2024parkin}. Its ubiquitin-like Ubl, RING0, RING1, IBR, REP/tether, and RING2 architecture couples folded zinc-binding domains to autoinhibitory linkers \citep{kumar2015parkin}. Of $8{,}756$ measured substitutions across $462$ profiled sites, $89.2\%$ of sites exhibit global--local separation. Raw \textsc{ESM-2} 650M with the wild-type interface reaches $0.503$ Spearman; MSA retrieval reaches $0.585$; \textsc{REM2} reaches $0.645$. Abundance-class AUC rises in parallel from $0.769$ to $0.815$ and $0.857$ (Figure~\ref{fig:viral_parkin}).

The rank shift has domain structure (Figure~\ref{fig:viral_parkin}d). \textsc{REM2} improves every folded domain, including RING0 from $0.425$ to $0.627$, IBR from $0.727$ to $0.786$, and RING2 from $0.427$ to $0.529$; REP/tether rises from $-0.083$ to $0.221$. The flexible Ubl--RING0 linker instead falls from $0.363$ to $0.176$, distinguishing the family-constrained catalytic core from a weakly conserved regime. N232A lies at the $0.4$th measured abundance percentile, yet Raw ranks it at the $88.8$th predicted percentile; MSA retrieval and \textsc{REM2} move it to predicted percentiles $58.0$ and $37.3$. For R396H, the predicted percentile moves in the opposite direction, from $39.7$ to $71.7$ and $78.0$, matching the measured $77.6$th percentile (Figure~\ref{fig:viral_parkin}e).

MSA retrieval produces its largest domain gains in RING0 ($+0.157$) and REP/tether ($+0.129$), while native-field calibration contributes a further $+0.065$ in RING2 and $+0.175$ in REP/tether; Ubl moves only $+0.032$ because its Raw ranking is already strong at $0.694$. The flexible linker, residues 77--143, is the sole negative response and is thereby set apart from the compact zinc-binding core.

\paragraph{Cross-protein region audit.}
The linker response is an instance of a general scope boundary rather than a Parkin-specific failure. For every ProteinGym assay, we split single substitutions by site confidence in the assay's AlphaFold2 model---flexible at pLDDT $<50$, ordered otherwise---and recompute within-region Spearman for Raw and \textsc{REM2} under the same \textsc{ESM-2} 650M wild-type readout, keeping subsets with at least $20$ measured variants at five or more sites. Ordered regions improve in $83\%$ of the $214$ qualifying assays (mean $\Delta=+0.056$; only $5.6\%$ degrade by more than $0.02$). Flexible regions expose the trade-off directly: across the $65$ qualifying assays, within-region ranking decreases for $57\%$, with mean $\Delta=-0.020$. Reliability weighting thus exchanges a small loss of resolution inside low-confidence regions for a larger gain across structured ones, and region-level degradations should be read against this scope condition rather than as isolated failures.

\FloatBarrier
\section{Readout, runtime, and scoring-interface audits}
\label{app:readout_runtime}

\subsection{Adaptive-stage and component audits}

\paragraph{Stages across backbones.}
Every configuration with complete staged fields is evaluated in the same order: Raw; entropy-weighted adaptive MSA mix; ungated background correction; coherence-gated correction; RSA shrinkage; and pLDDT shrinkage. The per-configuration view in Table~\ref{tab:staged_per_config} (ProteinGym), with the VenusMutHub and VenusViroHub counterparts in Tables~\ref{tab:staged_per_config_vmh}--\ref{tab:staged_per_config_viro}, shows that stage behavior is not ensemble-specific. Gated correction raises the aggregate mean on all three benchmarks, while RSA shrinkage improves all $59$ configurations on each benchmark.

\paragraph{Controlled \textsc{ProSST-2048} component isolation.}
Table~\ref{tab:ablation} applies every downstream operator to the same entropy-weighted MSA field and uses the complementary calibration coefficient $1-\alpha$. On the random 21-assay validation split, retrieval contributes $0.0201$ of the total $0.0276$ improvement; gating adds $0.0014$, and RSA and pLDDT shrinkage add a further $0.0061$, yielding $0.5849$ from a Raw value of $0.5573$. The held-out 196-assay ProteinGym test set follows the same order under official aggregation: retrieval contributes $0.0169$, gating $0.0040$, RSA shrinkage $0.0045$, and pLDDT shrinkage $0.0011$, producing a final score of $0.5337$ from $0.5072$. The ungated row is a gate intervention rather than a cumulative stage.

The coherence gate is evaluated independently of the component table. Relative to the same entropy-weighted adaptive mix, gated correction improves $59/59$, $56/59$, and $59/59$ configurations on ProteinGym, VenusMutHub, and VenusViroHub, respectively; relative to ungated correction, the corresponding counts are $59/59$, $23/59$, and $55/59$.

\paragraph{Fixed-share family-field sweep.}
\label{app:alpha_sweep}
Figure~\ref{fig:alpha_sweep} replaces the entropy-adaptive coefficient of Equation~\ref{eq:native_alpha} with a fixed family-field share $\alpha\in\{0.0,0.1,\dots,1.0\}$, evaluated on the same \textsc{ProSST}-2048 fields and the same 21-assay validation split as Table~\ref{tab:ablation}, under both the retrieval-only readout and the full pipeline. Three observations support the adaptive design. First, both endpoints degrade: $\alpha{=}0$ reduces to the Raw field and $\alpha{=}1$ collapses to $0.4093$, so neither channel suffices alone and the frozen backbone remains necessary even where family evidence is available. Second, performance forms a broad plateau over $\alpha\in[0.6,0.9]$ ($0.567$--$0.574$ retrieval-only), so the readout is not sensitive to the precise share within this range. Third, the per-protein entropy-adaptive coefficient reaches $0.5774$ retrieval-only and $0.5849$ with the full pipeline, matching or exceeding every fixed share without per-protein tuning. The plateau shape is unchanged when fold means are averaged over all ten folds of the partition ($0.531$--$0.537$ retrieval-only over $\alpha\in[0.6,0.9]$; $0.398$ at $\alpha{=}1$). Shaded intervals are 95\% assay-level bootstrap intervals over the 21 validation assays.

\begin{figure}[!ht]
\centering
\includegraphics[width=0.62\textwidth]{figures/fig_alpha_sweep.pdf}
\caption{Mean assay Spearman on the 21-assay validation split (\textsc{ProSST}-2048) under a fixed family-field share $\alpha$, for retrieval only and the full pipeline; bands are 95\% assay-bootstrap intervals. Dashed lines: the entropy-adaptive $\alpha$ of Equation~\ref{eq:native_alpha}. Both endpoints degrade, and the adaptive coefficient sits at or above the best fixed share.}
\label{fig:alpha_sweep}
\end{figure}

\input{table_gate_ablation.tex}

\input{table_staged_per_config.tex}

\input{table_staged_per_config_vmh.tex}

\input{table_staged_per_config_viro.tex}

\paragraph{Exposure-and-confidence shrinkage.}
\label{app:structure_extension}
$\lambda_i$ follows Equation~\ref{eq:lambda}: if RSA, pLDDT, or the sequence-to-structure map is unavailable, $\lambda_i=1$. Table~\ref{tab:ablation} isolates the two structural factors.

\subsection{Amortized score-field inference}

Figure~\ref{fig:runtime} measures post-forward processing after one shared \textsc{ProSST} field has been constructed per protein. Across 2.47 million ProteinGym variants, indexed lookup reduces mutation scoring from $114.56$ to $5.53$ seconds ($20.7\times$), while vectorized alignment construction reduces MSA-field construction from $3{,}026.75$ to $48.61$ seconds ($62.3\times$). Calibration and shrinkage add $O(L\times20)$ work. Model-resident execution, batched tokenization, vectorized field construction, and indexed lookup remove repeated work from mutation-wise loops.

\FloatBarrier
\subsection{Architecture-wide scoring-interface audit}
For a substitution $x_i\!\rightarrow\!a$, all three interfaces form the contrast $\Delta^R_{i,a}=R_{i,a}-R_{i,x_i}$ from Equation~\ref{eq:contrasts}, but they differ in the conditioning context used to produce $R$. A masked marginal uses $R^{\mathrm{mask}}_{i,a}=\log p_\theta(x_i=a \mid x_{\setminus i})$: the wild-type residue is masked and a full length-$L$ field requires $L$ encoder forward passes. A wild-type or unmasked field uses $R^{\mathrm{wt}}_{i,a}=\log p_\theta(x_i=a \mid x)$: a single forward pass scores all sites, but the wild-type residue remains in the input, so this is a ranking interface rather than a masked conditional. Teacher forcing uses $R^{\mathrm{tf}}_{i,a}=\log p_\theta(x_i=a \mid x_{<i})$ and includes structure when the backbone requires it (Equation~\ref{eq:readouts}). It supplies conditionals under a declared factorization, and one native-context field can be reused across variants. These definitions are held fixed within every Raw--\textsc{REM2} pair (Table~\ref{tab:score_interfaces}); \textsc{REM2} changes the readout field, not the conditioning convention.

\begin{table}[!ht]
\caption{Conditioning and cost of three scoring interfaces.}
\label{tab:score_interfaces}
\centering
\footnotesize
\setlength{\tabcolsep}{2pt}
\begin{tabular}{@{}>{\raggedright\arraybackslash}p{0.165\textwidth}>{\raggedright\arraybackslash}p{0.19\textwidth}>{\raggedright\arraybackslash}p{0.17\textwidth}>{\raggedright\arraybackslash}p{0.15\textwidth}>{\raggedright\arraybackslash}p{0.225\textwidth}@{}}
\toprule
Interface & Context at site $i$ & WT residue visible & Forward passes & Multi-mutant scoring \\
\midrule
Masked marginal & $x_{\setminus i}$ & No & $L$ per protein & Additive native field \\
Wild-type field & Full native $x$ & Yes & 1 per protein & Additive native field \\
Teacher forcing & Prefix $x_{<i}$ & \mbox{Yes (factorized)} & 1 per protein & \mbox{Additive unless rescored} \\
\bottomrule
\end{tabular}
\end{table}


\paragraph{Bidirectional encoders: masked versus wild-type fields.}
Across eight paired encoder checkpoints, both masked and wild-type interfaces improve on all three benchmarks, and the same global--local gain holds across conditioning conventions (Table~\ref{tab:mask_wt}).

\begin{table}[!ht]
\caption{Mean assay Spearman across eight paired checkpoints: \textsc{ESM-1b}, \textsc{ESM-1v}, \textsc{ESM-2} 8M/35M/150M/650M/3B, and \textsc{SaProt}. \textsc{REM2} uses entropy weighting. ProteinGym reports the assay mean, not official Average Spearman; VenusMutHub and VenusViroHub report flat assay macro-averages.}
\label{tab:mask_wt}
\centering
\small
\setlength{\tabcolsep}{7pt}
\begin{tabular}{lrrrrrr}
\toprule
Dataset & \multicolumn{3}{c}{Masked} & \multicolumn{3}{c}{Wild-type} \\
\cmidrule(lr){2-4}\cmidrule(lr){5-7}
& Raw & \textsc{REM2} & $\Delta$ & Raw & \textsc{REM2} & $\Delta$ \\
\midrule
ProteinGym & 0.389 & 0.468 & +0.079 & 0.395 & 0.458 & +0.062 \\
VenusMutHub & 0.157 & 0.218 & +0.061 & 0.142 & 0.187 & +0.045 \\
VenusViroHub & 0.102 & 0.265 & +0.163 & 0.133 & 0.300 & +0.168 \\
\bottomrule
\end{tabular}
\end{table}

\paragraph{Autoregressive and inverse-folding models: full-sequence likelihood versus a native field.}
For these architectures, full-sequence scoring teacher-forces each mutant and aggregates likelihood over the sequence, allowing the substitution to alter later conditionals. Native-field scoring instead teacher-forces the wild-type once and reads the mutant--wild-type contrast only at the mutated position. Table~\ref{tab:sequence_native_interfaces} pairs these routes on the same frozen checkpoints and ProteinGym assays.

\begin{table}[!ht]
\caption{ProteinGym interface comparison on 217 paired assays. Full-sequence values follow the official aggregate; $\Delta$ is native field minus full-sequence log-likelihood, and wins are paired at assay level.}
\label{tab:sequence_native_interfaces}
\centering
\small
\setlength{\tabcolsep}{6pt}
\begin{tabular}{llrrrr}
\toprule
Family & Model & Full sequence & Native field & $\Delta$ & Wins \\
\midrule
Inverse folding & \textsc{ESM-IF1} & 0.422 & 0.421 & $-0.001$ & 95/217 \\
Autoregressive & \textsc{ProGen2}-S & 0.336 & 0.288 & $-0.048$ & 47/217 \\
Autoregressive & \textsc{ProGen2}-M & 0.380 & 0.337 & $-0.043$ & 54/217 \\
Autoregressive & \textsc{ProGen2}-B & 0.378 & 0.330 & $-0.048$ & 58/217 \\
Autoregressive & \textsc{ProGen2}-L & 0.380 & 0.327 & $-0.053$ & 53/217 \\
Autoregressive & \textsc{ProGen2}-XL & 0.391 & 0.351 & $-0.040$ & 67/217 \\
\bottomrule
\end{tabular}
\end{table}

The two routes induce almost the same \textsc{ESM-IF1} ranking (assay-level Spearman $0.966$), consistent with structure conditioning localizing the likelihood change. For \textsc{ProGen2}, however, full-sequence likelihood retains useful downstream context and exceeds the native field at every scale. The reusable field is therefore an architecture-compatible interface for \textsc{REM2}, not an assertion that discarding all contextual propagation must improve the backbone score.

\FloatBarrier
\subsection{\textsc{ProteinMPNN} reverses the autoregressive trend}

\paragraph{Official score versus teacher forcing.}
Unlike the stable \textsc{ESM-IF1} comparison and the full-sequence advantage of \textsc{ProGen2}, official \textsc{ProteinMPNN} scoring samples a decoding order while rebuilding each mutant sequence, making the score order-dependent. Our native teacher-forcing route evaluates one conditional amino-acid field on the native sequence and applies mutant-minus-wild-type lookup (Table~\ref{tab:pmpnn_interfaces}). On identical AlphaFold2 structures, this field raises mean assay Spearman from $0.271$--$0.290$ to $0.376$--$0.401$ across eight checkpoints while reducing wall-clock cost by $219$--$601\times$ (Figure~\ref{fig:tf}a); at v\_48\_020, $204$ of $217$ assays improve (Figure~\ref{fig:tf}b). The sign reversal localizes the gain to removal of decoding-path variability rather than teacher forcing alone. Conditional rescoring then exposes a non-additive component for the following epistasis diagnostic.

\begin{table}[!ht]
\caption{\textsc{ProteinMPNN} official random-order versus native teacher-forcing (TF) scoring.}
\label{tab:pmpnn_interfaces}
\centering
\footnotesize
\setlength{\tabcolsep}{2.5pt}
\begin{tabular}{@{}lllll@{}}
\toprule
Route & Context & Evaluations & Order & Multi-mutant score \\
\midrule
Official random order & Mutant + structure & Per mutant & Yes & May be non-additive \\
Native TF & Native + structure & Per protein & No & Additive \\
Conditional TF & Mutant background + structure & Per background & No & Non-additive \\
\bottomrule
\end{tabular}
\end{table}

\begin{figure}[!ht]
\begin{center}
\includegraphics[width=0.95\textwidth]{figures/fig_teacherforce.pdf}
\end{center}
\caption{\textsc{ProteinMPNN} teacher forcing versus official random-order scoring on AlphaFold2 structures. \textbf{(a)} Mean assay Spearman and wall-clock cost across eight checkpoints. \textbf{(b)} Per-assay Spearman at v\_48\_020.}
\label{fig:tf}
\end{figure}

\FloatBarrier
\subsection{Pairwise epistasis diagnostic}

\paragraph{\textsc{ProteinMPNN} epistasis evaluation.}
We evaluate epistasis on all 535{,}917 complete wild-type/single/single/double quartets in the GB1 binding landscape of \citet{olson2014epistasis}. Teacher-forced variant scores are additive by construction, with maximum numerical deviation $4.8\times10^{-7}$, so we also rescore mutation A on the B background and vice versa. Table~\ref{tab:pmpnn_epistasis} separates main effects from interactions: teacher forcing is substantially stronger for single- and double-mutant fitness, while epistasis correlations remain near zero for both the official random-order score, $0.005$ with interval $[-0.002,0.013]$, and conditional teacher forcing, $0.012$ with interval $[-0.005,0.029]$.

\begin{table}[H]
\caption{\textsc{ProteinMPNN} v\_48\_020 on Olson GB1; $\rho_\epsilon$ is Spearman correlation with measured pairwise epistasis.}
\label{tab:pmpnn_epistasis}
\begin{center}
\scriptsize
\setlength{\tabcolsep}{3pt}
\resizebox{\textwidth}{!}{%
\begin{tabular}{lrrrrrrr}
\toprule
Subset & Quartets & Single TF & Single official & Double TF & Double official & $\rho_\epsilon$ official & $\rho_\epsilon$ conditional TF \\
\midrule
All & 535{,}917 & 0.306 & 0.105 & 0.284 & 0.165 & 0.005 & 0.012 \\
Non-floor & 400{,}118 & 0.312 & 0.091 & 0.292 & 0.152 & 0.012 & 0.022 \\
\bottomrule
\end{tabular}}
\end{center}
\end{table}

The focused GB1 diagnostic separates representation from readout: the official and conditional routes recover little pairwise-interaction signal, while the lower-cost additive field retains substantially more main-effect signal under the same frozen backbone and structure inputs.

\FloatBarrier
\section{Extended ProteinGym metrics and stratified analyses}
\label{app:extended_metrics}

Per-configuration NDCG, AUC, MCC, and top-variant recall are in Tables~\ref{tab:pg_metrics_family_1}--\ref{tab:viro_metrics_family_4}. These tables are the configuration-level source of Figure~\ref{fig:scatter}: all $295$ model--metric pairs improve on ProteinGym and VenusViroHub, and $293$ of $295$ improve on VenusMutHub. Table~\ref{tab:stratified_taxa_msa} then asks whether the \textsc{ProSST}-ensemble Raw--\textsc{REM2} gain, already visible as the \textsc{VenusREM2} row of Table~\ref{tab:leaderboard_full_taxon_msa}, is confined to one taxonomic or MSA-depth bin.

\input{table_stratified_taxa_msa.tex}

The \textsc{ProSST}-ensemble gain is not confined to one bin: virus assays improve by $0.064$ from $0.476$ to $0.540$, and the official low-depth bin gains $0.055$ from $0.498$ to $0.553$; even the high-depth bin increases by $0.011$, from $0.612$ to $0.623$. The five-metric family blocks follow, grouped by benchmark and ordered from structure-conditioned (structure-aware and inverse-folding) to sequence-only (masked and autoregressive) models.

\input{table_model_metrics.tex}

\clearpage
\section{Cross-benchmark analyses of native backgrounds and family evidence}
\label{app:replications}

\subsection{ProteinGym}
\label{app:proteingym_replication}
Figure~\ref{fig:proteingym_analysis} gives the full-benchmark aggregate counterpart of the native-background and family-gain analyses using the same 37 canonical checkpoints and all 217 ProteinGym assays.

\begin{figure}[H]
\begin{center}
\includegraphics[width=0.94\textwidth]{figures/figA1_proteingym_analysis.pdf}
\end{center}
\caption{ProteinGym replication of the native-background and motivation analyses. \textbf{(a)} Checkpoint amino-acid backgrounds and their agreement with global UniProt composition. \textbf{(b)} Assay-matched native background versus MSA family composition; the central 99\% of amino-acid pairs is shown, while correlations use all $217\times20$ pairs. \textbf{(c)} Average-Spearman Raw--\textsc{REM2} performance by backbone family.}
\label{fig:proteingym_analysis}
\end{figure}

Checkpoint backgrounds have median pairwise Pearson $r=0.995$ and consensus--corpus correlation $r=0.993$ (Figure~\ref{fig:proteingym_analysis}a). Across 217 assay-matched alignments, native and family composition agree at median $r=0.965$ and pooled $r=0.758$ (Figure~\ref{fig:proteingym_analysis}b). \textsc{REM2} improves all $59$ configurations, with mean gains of $0.037$ for structure-aware models ($0.456$ to $0.493$), $0.066$ for masked sequence models ($0.374$ to $0.440$), $0.063$ for inverse-folding models ($0.382$ to $0.445$), and $0.106$ for autoregressive models ($0.302$ to $0.408$) (Figure~\ref{fig:proteingym_analysis}c).

\clearpage
\subsection{VenusMutHub}
\label{app:vmh}
Figure~\ref{fig:vmh_analysis} tests aggregate native-background consistency and family-level gains on VenusMutHub using 37 checkpoints and 905 query-matched alignments. Figure~\ref{fig:motivation}b additionally applies the site-level construction to complete cached \textsc{ESM-2} 650M WT and MSA fields after duplicate assays are collapsed by WT sequence.

\begin{figure}[H]
\begin{center}
\includegraphics[width=0.94\textwidth]{figures/figA2_vmh_analysis.pdf}
\end{center}
\caption{VenusMutHub analysis of native backgrounds and family evidence. \textbf{(a)} Checkpoint backgrounds versus global UniProt composition. \textbf{(b)} Query-matched native background versus MSA composition; the central 99\% is shown, while correlations use all $905\times20$ pairs. \textbf{(c)} Assay-macro Raw--\textsc{REM2} Spearman by backbone family.}
\label{fig:vmh_analysis}
\end{figure}

The shared checkpoint background reappears with median pairwise Pearson $r=0.997$ and corpus correlation $r=0.990$ (Figure~\ref{fig:vmh_analysis}a). Across 905 query-matched alignments, native and family composition agree at median $r=0.964$ and pooled $r=0.863$ (Figure~\ref{fig:vmh_analysis}b). At site resolution, $148{,}983$ complete profiles from $527$ unique WT sequences have median correlations of $0.249$ with corpus abundance and $0.229$ with the protein-level native background; only $24.8\%$ align with both at $r\geq0.3$ (Figure~\ref{fig:motivation}b). \textsc{REM2} improves all $59$ configurations, with mean gains of $0.052$ for structure-aware models ($0.176$ to $0.228$), $0.055$ for masked sequence models ($0.145$ to $0.200$), $0.046$ for inverse-folding models ($0.210$ to $0.256$), and $0.080$ for autoregressive models ($0.102$ to $0.182$) (Figure~\ref{fig:vmh_analysis}c). Across all 905 assays, MSA retrieval contributes $0.058$ to the \textsc{ProSST} ensemble, while full \textsc{REM2} reaches $0.258$, a total gain of $0.068$ over the Raw value of $0.190$.

\clearpage
\subsection{VenusViroHub}
\label{app:virogym_replication}
Figure~\ref{fig:virogym_analysis} tests the aggregate background and repeats the site-level separation analysis under an independent virus-focused distribution using the same 37 canonical checkpoints, 89 assays, and 52 unique wild-type sequences.

\begin{figure}[H]
\begin{center}
\includegraphics[width=0.94\textwidth]{figures/figA3_virogym_analysis.pdf}
\end{center}
\caption{VenusViroHub analysis of native backgrounds and family evidence. \textbf{(a)} Checkpoint backgrounds versus global UniProt composition. \textbf{(b)} Wild-type native background versus aggregate MSA composition. \textbf{(c)} Independent site-level replication of Figure~\ref{fig:motivation}b. \textbf{(d)} Hierarchically aggregated family-level Raw--\textsc{REM2} Spearman.}
\label{fig:virogym_analysis}
\end{figure}

The frozen native backgrounds remain highly concordant across checkpoints, with median pairwise Pearson $r=0.990$, while the consensus retains a corpus correlation of $r=0.885$ (Figure~\ref{fig:virogym_analysis}a). Within the 52 wild-type sequences, native background and MSA family composition agree at median $r=0.976$; the lower pooled correlation, $r=0.621$, separates backbone-specific composition from a single global abundance profile (Figure~\ref{fig:virogym_analysis}b). After repeated phenotype assays are collapsed by wild-type sequence, the same \textsc{ESM-2} 650M wild-type-interface construction yields $29{,}416$ complete site profiles. Their median correlations are $0.273$ with corpus abundance and $0.332$ with the protein-level native background; only $38.1\%$ align with both at $r\geq0.3$ (Figure~\ref{fig:virogym_analysis}c).

The benchmark-level outcome is consistent with this complementary signal. \textsc{REM2} improves all $59$ configurations: masked sequence models rise from $0.106$ to $0.246$ and autoregressive models from $0.121$ to $0.228$, while structure-aware and inverse-folding families improve from $0.240$ to $0.301$ and from $0.288$ to $0.325$ (Figure~\ref{fig:virogym_analysis}d). Together with VenusMutHub, these analyses show that the shared native background and the corrective family signal persist across both high-precision small-data assays and an independent viral benchmark.

\end{document}
